\section{Results}
\label{sec:results}

\subsection{Main Result: Verbosity Coefficient Reversal (H-M2)}

The primary evidence for annotation drift is the reversal of verbosity preference sign across annotation strata. Table~\ref{tab:coefficients} reports the round-stratified logistic regression coefficients with 95\% bootstrap confidence intervals (2,000 stratified resamples on 160,800 HH-RLHF preference comparisons).

\begin{table}[h]
\centering
\caption{Stylistic preference coefficients across annotation rounds (H-M2). Bold indicates CI non-overlap.}
\label{tab:coefficients}
\begin{tabular}{lllllll}
\toprule
Feature & Early $\beta$ & 95\% CI & Late $\beta$ & 95\% CI & $\Delta$ & Non-Overlap \\
\midrule
$\beta_L$ (verbosity) & $-0.025$ & $[-0.043, -0.006]$ & $+0.056$ & $[+0.043, +0.068]$ & \textbf{+0.080} & \textbf{YES} \\
$\beta_H$ (hedging) & $-0.029$ & $[-0.048, -0.011]$ & $-0.008$ & $[-0.024, +0.007]$ & $+0.021$ & no \\
$\beta_S$ (structure) & $-0.002$ & $[-0.021, +0.010]$ & $+0.010$ & $[+0.004, +0.016]$ & $+0.012$ & no \\
$\beta_Q$ (quality) & --- & --- & $-0.017$ & --- & --- & stable \\
\bottomrule
\end{tabular}
\end{table}

\begin{figure}[h]
\centering
\includegraphics[width=0.7\linewidth]{figures/fig1_coefficient_comparison.png}
\caption{Stylistic preference coefficients across annotation rounds with 95\% bootstrap confidence intervals. Early-round (round 1) and late-round (round 3) logistic regression coefficients for verbosity ($\beta_L$), hedging ($\beta_H$), and structured reasoning ($\beta_S$), estimated on HH-RLHF (53,600 rows per stratum, 2,000 bootstrap resamples). $\beta_L$ CIs are non-overlapping ($\Delta = +0.080$); $\beta_H$ and $\beta_S$ CIs overlap.}
\label{fig:coefficients}
\end{figure}

Three key observations: (1) \emph{Verbosity preference reverses direction.} The 95\% bootstrap CIs are non-overlapping --- a direction-level shift of $\Delta\beta_L = +0.080$. Early annotators penalized verbose responses; later annotators rewarded them. (2) \emph{Directional consistency across all features.} All three stylistic deltas are positive (sign\_consistent = true). (3) \emph{Quality covariate stable.} $\beta_Q = -0.017$, well within the $|\beta_Q| < 0.2$ threshold, confirming the verbosity shift is not quality recalibration.

\begin{figure}[h]
\centering
\includegraphics[width=0.7\linewidth]{figures/fig3_feature_stability_rounds.png}
\caption{Feature coefficient stability across annotation rounds. Trajectory of $\beta_L$, $\beta_H$, and $\beta_S$ point estimates from round 1 to round 3. All three features trend positive (sign\_consistent = true), with $\beta_L$ showing the largest and only CI-non-overlapping shift.}
\label{fig:stability}
\end{figure}

\subsection{AI-Typicality Geometric Projection (H-M1)}

The between-group regression on WebGPT (between-group tercile design) yields:

\begin{equation}
\beta_{\text{exposure}} = 0.041 \quad (p = 2.05 \times 10^{-5}; \text{ tercile } F = 82.92, p \approx 1.4 \times 10^{-36})
\end{equation}

\begin{figure}[h]
\centering
\includegraphics[width=0.7\linewidth]{figures/dose_response.png}
\caption{AI-typicality projection scores across annotator terciles (dose-response). Between-group regression of AI-typicality projection score on exposure tercile (WebGPT, $n=19{,}578$). Higher terciles show stronger alignment with the AI-typicality direction. $\beta_\text{exposure} = 0.041$, $p = 2.05\times10^{-5}$.}
\label{fig:dose_response}
\end{figure}

\textbf{Discriminant validity.} The placebo permutation test (200 random vectors) yields $p = 0.48$, compared to observed $p = 2.05\times10^{-5}$. The projection signal is specific to the AI-typicality direction --- not a general embedding artifact.

\begin{figure}[h]
\centering
\includegraphics[width=0.7\linewidth]{figures/discriminant_validity.png}
\caption{Discriminant validity --- AI-typicality vs.\ placebo projection. The observed $\beta_\text{exposure} = 0.041$ lies far outside the placebo null distribution from 200 permutations (empirical placebo $p = 0.48$), confirming signal specificity.}
\label{fig:discriminant}
\end{figure}

\subsection{Null Results as Methodological Contribution (H-E1)}

The round$\times$ambiguity interaction returned $p = 1.0$ and zero high-ambiguity samples were detected. This is not falsification --- HH-RLHF lacks per-prompt disagreement labels, making the interaction test architecturally impossible. H-E1 confirms $\beta_L$ nominal Bonferroni significance ($p = 0.000$) and feature orthogonality (VIF $< 1.03$).

\subsection{Cross-Experiment Summary}

\begin{table}[h]
\centering
\caption{Cross-experiment results summary.}
\label{tab:summary}
\begin{tabular}{llllll}
\toprule
Experiment & Type & Gate & Result & Primary Metric & Value \\
\midrule
H-E1 & Existence & MUST\_WORK & PASS & Interaction $p$ & 1.0 (data limit) \\
H-M1 & Mechanism & MUST\_WORK & PASS & $\beta_\text{exposure}$ & 0.041, $p=2.05\times10^{-5}$ \\
H-M2 & Mechanism & SHOULD\_WORK & PARTIAL & $n\_\text{directional}$ & 1/3; sign\_consistent=true \\
\bottomrule
\end{tabular}
\end{table}

\subsection{Informative Null Results}

\begin{table}[h]
\centering
\caption{Informative null results and data requirements.}
\label{tab:nullresults}
\begin{tabular}{lll}
\toprule
Null Result & Reason & Required Data \\
\midrule
Interaction $p = 1.0$ & No per-prompt disagreement labels & Multi-annotator Fleiss $\kappa$ per prompt \\
Effect size $< 0.1$ SD & No within-annotator worker IDs & Genuine session tracking ($\geq$3 sessions/worker) \\
$\beta_H$, $\beta_S$ subthreshold & Index proxy + topic imbalance & Topic-balanced strata + larger round depth \\
\bottomrule
\end{tabular}
\end{table}

\begin{figure}[h]
\centering
\includegraphics[width=0.7\linewidth]{figures/fig5_topic_balance.png}
\caption{Topic distribution imbalance across HH-RLHF strata. Chi-square test for topic uniformity ($p = 4\times10^{-275}$), indicating extreme non-uniform topic distribution across annotation strata.}
\label{fig:topic_balance}
\end{figure}
